%% =============================================================================
%% NRL Report LaTeX Template for PaperOrchestra / Multiagent Authoring Systems
%% =============================================================================
%%
%% This template is designed to be programmatically filled by a multiagent
%% framework (e.g., PaperOrchestra) to produce a complete NRL Memorandum Report
%% or NRL Formal Report.
%%
%% DOCUMENT CLASS: nrlreport.cls (must be in the same directory or installed)
%% REQUIRED FILES IN WORKING DIRECTORY:
%%   - nrlreport.cls          (the NRL report document class)
%%   - bibunits.sty           (bibliography units support for appendix bibs)
%%   - doi.sty                (DOI hyperlink support)
%%   - NRL_logo_1C.pdf        (NRL logo for the title page)
%%   - nrlunsrtabbrvnat.bst   (default bibliography style -- citation order, abbreviated names)
%%   - References.bib         (your BibTeX database)
%%
%% AVAILABLE BIBLIOGRAPHY STYLES (choose one via \bibliographystyle):
%%   - nrlunsrtabbrvnat.bst   (DEFAULT: citation order, abbreviated names)
%%   - nrlunsrtnat.bst        (citation order, full names)
%%   - nrlabbrvnat.bst        (alphabetical order, abbreviated names)
%%   - nrlplainnat.bst        (alphabetical order, full names)
%%
%% COMPILATION SEQUENCE:
%%   pdflatex template        % first pass: compile lists of citations
%%   bibtex template          % extract citations for main paper
%%   bibtex bu1               % extract citations for first appendix (if any)
%%   bibtex bu2               % extract citations for second appendix (if any)
%%   pdflatex template        % second pass: typeset the bibliographies
%%   pdflatex template        % third pass: number the citations and finalize
%%
%% ALTERNATIVE COMPILERS: lualatex and xelatex are also supported.
%%   For unicode support, use: \documentclass[unicode]{nrlreport}
%%
%% CLASS OPTIONS (all optional, defaults in parentheses):
%%   Font size: 10pt, (11pt), 12pt
%%   Spacing:   (singlespace), doublespace
%%   Chapter start: (openright), openany
%%   Equation alignment: (centered), fleqn (flush-left)
%%   Draft mode example: \documentclass[12pt,doublespace,openany,fleqn]{nrlreport}
%%
%% =============================================================================

\documentclass{nrlreport} % final report mode
%\documentclass[12pt,doublespace,openany,fleqn]{nrlreport} % nice draft mode

%% =============================================================================
%% PREAMBLE: All \nrl commands below MUST appear before \maketitle.
%% If you do not know a value, omit the command entirely -- all have defaults.
%% =============================================================================

%% -- Bibliography database for appendix bibliographies (bibunits) -------------
%% Point this to your .bib file (without the .bib extension).
\defaultbibliography{References}

%% -- Report Number ------------------------------------------------------------
%% Format: \newreportnumber{branch-code}{FR or MR}{4-digit-year}{sequence-number}
%%   FR = Formal Report, MR = Memorandum Report
%% PLACEHOLDER: Replace with actual values once assigned by TIS.
\newreportnumber{XXXX}{MR}{YYYY}{XX}

%% -- Security Classification --------------------------------------------------
%% \nrlmark sets the classification banner on the top and bottom of every page.
%% Leave empty {} for unclassified. Examples: {UNCLASSIFIED}, {CUI}, {SECRET}
%\nrlmark{}

%% Bold distribution mark at the bottom of the title page.
%% Default is "NOT APPROVED FOR DISTRIBUTION". Set to {} to turn off.
\nrldistributionmark{}

%% Security statements for the title page (uncomment and fill as needed):
%\nrlwarning{}           % Predefined: {NF}, {EX}, {TTCP}, or {custom warning}
%\nrldistribution{A}     % Predefined: {A}, {B}, {C}, {D}, {E}, {F}, or {custom}
%\nrldissemination{}     % Optional clarification on dissemination

%% Classification authority (triggers CUI caveat box on title page):
%\nrlclassifiedby{}
%\nrlderivedfrom{}
%\nrldeclassifyon{}

%% CUI information box (uncomment and fill as needed):
%\nrlCUIcontrolledbyComponent{}
%\nrlCUIcontrolledbyOffice{}
%\nrlCUIcategories{}
%\nrlCUIdistribution{}
%\nrlCUIdissemenation{}
%\nrlCUIpoc{}
%\nrlCUIrefpoc{}
%\nrlCUIreference{}

%% -- Report Documentation Page Fields -----------------------------------------
%% These populate the standard SF298-style documentation page generated by \nrlabstract.
%\nrlreporttype{NRL Memorandum Report}  % or: NRL Formal Report
%\nrlstartdate{}            % Project start date
%\nrlenddate{}              % Project end date
%\nrlauthorlist{}           % Comma-separated author list for documentation page

%\nrlcontractnumber{}
%\nrlgrantnumber{}
%\nrlprogramelement{}
%\nrlprojectnumber{}
%\nrltasknumber{}
%\nrlworkunitnumber{}

%\nrlorganizationaddress{}
%\nrlsponsoraddress{}
%\nrlsponsoracronym{}       % e.g., ONR
%\nrlsponsorreportnumber{}

%\nrlsupplementarynotes{}
%\nrlsubjectterms{}         % Keywords / subject terms

%\nrlreportclass{}          % Classification of the report
%\nrlabstractclass{}        % Classification of the abstract
%\nrlthispageclass{}        % Classification of the documentation page
%\nrlabstractlimitation{}
%\nrlresponsibleperson{}
%\nrltelephonenumber{}

%% -- Additional Packages (add any extra packages your content needs) ----------
%% The nrlreport class already loads: amsmath, graphicx, hyperref, natbib, etc.
% \usepackage{algorithm2e}
% \usepackage{listings}
% \usepackage{booktabs}
% \usepackage{multirow}
% \usepackage{subcaption}

%% =============================================================================
%% DOCUMENT BODY
%% =============================================================================

\begin{document}

%% -- Title and Author Information ---------------------------------------------
%% NOTE: \title, \author, \nrlheadertitle, \nrlheaderauthor, and \nrlbranch
%%       MUST appear before \maketitle. They may appear in the preamble or here.

\title{PLACEHOLDER REPORT TITLE}

%% Short title for the running header on odd-numbered pages:
\nrlheadertitle{PLACEHOLDER SHORT TITLE}

%% Optional: if the title page title needs different formatting than \title:
%\nrlpagetitle{}

%% Author for the running header on even-numbered pages (typically first author):
\nrlheaderauthor{PLACEHOLDER FIRST AUTHOR}

%% Full author list for the title page (all authors in the same branch together):
\author{PLACEHOLDER AUTHOR ONE, PLACEHOLDER AUTHOR TWO}

%% Branch and division for the primary author group:
\nrlbranch{PLACEHOLDER Branch Name\\PLACEHOLDER Division Name}

%% Additional author groups from other branches (up to 5 total):
%\nrlauthortwo{Author Three}
%\nrlbranchtwo{Second Branch Name\\Second Division Name}
%\nrlauthorthree{Author Four}
%\nrlbranchthree{Third Branch Name\\Third Division Name}
%\nrlauthorfour{}
%\nrlbranchfour{}
%\nrlauthorfive{}
%\nrlbranchfive{}

%% -- Dates --------------------------------------------------------------------
%% Standard LaTeX \date{} is IGNORED by this class. Use these instead:
\nrltitlepagedate{PLACEHOLDER Month DD, YYYY}    % Date on the title page
%\nrlreportdate{DD-MM-YYYY}                       % Date on the documentation page
%\nrlapprovaldate{Month DD, YYYY}                 % Date in the footnote on page 1

%% -- Generate the Title Page --------------------------------------------------
\maketitle

%% -- Abstract (Report Documentation Page) -------------------------------------
%% \nrlabstract generates the SF298 Report Documentation Page.
%% It MUST appear after \maketitle.
\nrlabstract{%
PLACEHOLDER: Provide a concise abstract (typically 200 words or fewer) summarizing
the purpose, methodology, key findings, and conclusions of this report.
}

%% -- Front Matter -------------------------------------------------------------
\tableofcontents    % Required by TIS
\listoffigures      % Optional
\listoftables       % Optional

\begin{executivesummary}
PLACEHOLDER: Provide an extended abstract / executive summary of the report.
This should give readers a self-contained overview of the work, including
motivation, approach, principal results, and significance.
\end{executivesummary}


%% =============================================================================
%% MAIN BODY CHAPTERS
%% =============================================================================
%% IMPORTANT: The top-level sectioning command is \chapter (NOT \section).
%% Hierarchy: \chapter > \section > \subsection > \subsubsection > \paragraph
%%
%% NRL cross-reference macros (use these for uniform formatting):
%%   \Secn{label} / \secn{label}    -- "Section X" / "Section X"
%%   \Eqn{label}  / \eqn{label}     -- "Equation (X)" / "Eq.(X)"
%%   \Fig{label}  / \fig{label}     -- "Figure X" / "Fig.X"
%%   \Tbl{label}  / \tbl{label}     -- "Table X" / "Table X"
%%   \App{label}  / \app{label}     -- "Appendix X" / "Appendix X"
%%   Range variants: \Eqns, \Figs, \Secs (two args for start--end)
%%   And variants:   \Eqnand, \Figand, \Secand (two args)
%%
%% Physical units macro: \units{\frac{m}{s}} produces upright-roman units.
%%
%% Paragraph classification markings (for classified documents):
%%   \u = (U), \f = (FOUO), \cui = (CUI), \c = (C), \s = (S), etc.
%% =============================================================================

\chapter{INTRODUCTION}
\label{chap:introduction}

PLACEHOLDER: Introduce the problem domain, motivation, and objectives of this work.
State the research questions or hypotheses. Provide an overview of the report structure.

%% Example citation: \cite{key} produces a numbered reference.
%% Example cross-ref: As discussed in \secn{chap:methods}, ...


\chapter{BACKGROUND AND RELATED WORK}
\label{chap:background}

PLACEHOLDER: Review relevant prior work, theoretical foundations, and context
necessary for understanding the contributions of this report.

\section{Theoretical Framework}
\label{sec:theory}

PLACEHOLDER: Present the theoretical basis for this work.

%% Example equation:
%% \begin{equation}
%%   E = mc^2
%%   \label{eq:example}  % label should be the LAST line in equation env
%% \end{equation}
%% Referenced as: \Eqn{eq:example} or \eqn{eq:example}.

\section{Prior Work}
\label{sec:prior_work}

PLACEHOLDER: Discuss prior work and how this report builds upon or differs from it.


\chapter{METHODOLOGY}
\label{chap:methods}

PLACEHOLDER: Describe the methods, algorithms, experimental design, or analytical
approach used in this work. Provide enough detail for reproducibility.

\section{Problem Formulation}
\label{sec:formulation}

PLACEHOLDER: Formally define the problem being addressed.

\section{Approach}
\label{sec:approach}

PLACEHOLDER: Describe the proposed approach in detail.

%% Example figure (standard LaTeX):
%% \begin{figure}
%%   \includegraphics[width=0.8\columnwidth]{filename}
%%   \caption{Description of the figure.}
%%   \label{fig:example}
%% \end{figure}
%%
%% Example figure (NRL environment with optional classification box):
%% \begin{nrlfigure}{UNCLASSIFIED}{Short caption}{Long caption.\label{fig:example2}}
%%   \includegraphics[width=0.8\columnwidth]{filename}
%% \end{nrlfigure}
%%
%% Referenced as: \Fig{fig:example} or \fig{fig:example}.


\chapter{RESULTS}
\label{chap:results}

PLACEHOLDER: Present the results of the work. Use figures, tables, and equations
as appropriate to convey findings clearly.

\section{Experimental Setup}
\label{sec:setup}

PLACEHOLDER: Describe the experimental conditions, datasets, parameters, etc.

\section{Analysis}
\label{sec:analysis}

PLACEHOLDER: Analyze the results. Compare against baselines, prior work,
or theoretical predictions.

%% Example table (standard LaTeX):
%% \begin{table}
%%   \caption{Description of the table.}
%%   \label{tab:example}
%%   \begin{tabular}{|l|c|r|}
%%     \hline
%%     Column 1 & Column 2 & Column 3 \\
%%     \hline
%%     data & data & data \\
%%     \hline
%%   \end{tabular}
%% \end{table}
%%
%% Example table (NRL environment with optional classification box):
%% \begin{nrltable}{UNCLASSIFIED}{Short caption}{Long caption.\label{tab:example2}}
%%   \begin{tabular}{|l|c|r|}
%%     \hline
%%     Column 1 & Column 2 & Column 3 \\
%%     \hline
%%   \end{tabular}
%% \end{nrltable}
%%
%% Referenced as: \Tbl{tab:example} or \tbl{tab:example}.


\chapter{DISCUSSION}
\label{chap:discussion}

PLACEHOLDER: Interpret the results in the broader context. Discuss implications,
limitations, and potential applications of the findings.


\chapter{SUMMARY AND CONCLUSIONS}
\label{chap:conclusions}

PLACEHOLDER: Summarize the key contributions and findings. State conclusions
and recommend directions for future work.


%% =============================================================================
%% ACKNOWLEDGMENTS (optional)
%% =============================================================================

\begin{acknowledgments}
PLACEHOLDER: Acknowledge funding sources, collaborators, and supporting organizations.
\end{acknowledgments}


%% =============================================================================
%% BIBLIOGRAPHY
%% =============================================================================
%% Point this to your .bib file (without the .bib extension).
%% The bibliography style is set automatically by the nrlreport class
%% (default: nrlunsrtabbrvnat). To override, add \bibliographystyle{style}
%% before \bibliography.

\bibliography{References}


%% =============================================================================
%% APPENDICES (optional)
%% =============================================================================
%% Each appendix gets its own \chapter and its own bibliography via bibunits.
%% Use Initial Caps in appendix chapter names (required for correct TOC formatting).
%% The \defaultbibliography command in the preamble points bibunits to your .bib file.
%%
%% COMPILATION NOTE: Each appendix bibliography requires a separate bibtex run:
%%   bibtex bu1   (for the first appendix)
%%   bibtex bu2   (for the second appendix)
%%   etc.

\appendix

\chapter{Supplementary Material}
\label{app:supplementary}

\begin{bibunit}

PLACEHOLDER: Include supplementary derivations, additional data tables,
extended results, or other material that supports but is not essential
to the main narrative.

%% Uncomment the following line to produce the appendix bibliography:
%\putbib

\end{bibunit}


%% Additional appendices follow the same pattern:
%
% \chapter{Additional Appendix Title}
% \label{app:additional}
%
% \begin{bibunit}
% Content here...
% \putbib
% \end{bibunit}


\end{document}
